\section{Koherencja Mac Lane'a i funktory monoidalne}\label{cat:sec:coherence}

Pracujemy z kategorią monoidalną w sensie definicji~\ref{cat:def:monoidal}. Najpierw wyprowadzimy tożsamości jednostki, następnie udowodnimy koherencję i skonstruujemy kategorię ściśle monoidalną równoważną wyjściowej.

\subsection{Dodatkowe tożsamości jednostki}
\begin{lemat}[Tożsamości jednostki]\label{cat:lem:unit-coherence}
W każdej kategorii monoidalnej zachodzą
\begin{align}
\lambda_{A\otimes B}\alpha_{I,A,B}&=\lambda_A\otimes\id_B,\label{cat:eq:left-unit-coherence}\\
\rho_{A\otimes B}&=(\id_A\otimes\rho_B)\alpha_{A,B,I},\label{cat:eq:right-unit-coherence}\\
\lambda_I&=\rho_I.\label{cat:eq:unitors-equal}
\end{align}
\end{lemat}
\begin{proof}
W dowodzie opuszczamy znak składania, zachowując kolejność od prawej do lewej. Funktory $I\otimes-$ oraz $-\otimes I$ są wierne, ponieważ są naturalnie izomorficzne z funktorem identycznościowym. Dokładniej, równość $I\otimes f=I\otimes g$ po złożeniu z $\lambda$ i wykorzystaniu naturalności daje $f\lambda=g\lambda$, skąd $f=g$.

Ustalmy $X=I$ (poniższe rachunki działają również dla dowolnego $X$). Pięciokąt dla $X,I,A,B$ daje
\[
\begin{aligned}
&\alpha_{X,I,A\otimes B}\alpha_{X\otimes I,A,B}\\
&=(\id_X\otimes\alpha_{I,A,B})
\alpha_{X,I\otimes A,B}(\alpha_{X,I,A}\otimes\id_B).
\end{aligned}
\]
Składamy obie strony z $\id_X\otimes\lambda_{A\otimes B}$. Lewa strona, dzięki trójkątowi i naturalności $\alpha$, staje się
\[
\begin{aligned}
&(\rho_X\otimes\id_{A\otimes B})\alpha_{X\otimes I,A,B}\\
&=\alpha_{X,A,B}((\rho_X\otimes\id_A)\otimes\id_B)\\
&=\alpha_{X,A,B}((\id_X\otimes\lambda_A)\otimes\id_B)
(\alpha_{X,I,A}\otimes\id_B)\\
&=(\id_X\otimes(\lambda_A\otimes\id_B))
\alpha_{X,I\otimes A,B}(\alpha_{X,I,A}\otimes\id_B).
\end{aligned}
\]
Druga równość to naturalność, trzecia wynika z trójkąta (po uwzględnieniu bifunktorialności), a ostatnia z naturalności $\alpha$. Prawa strona pięciokąta po złożeniu to
\[
(\id_X\otimes(\lambda_{A\otimes B}\alpha_{I,A,B}))
\alpha_{X,I\otimes A,B}(\alpha_{X,I,A}\otimes\id_B).
\]
Skracamy wspólne izomorfizmy po prawej, następnie korzystamy z wierności $I\otimes-$, otrzymując \eqref{cat:eq:left-unit-coherence}.

Aby otrzymać \eqref{cat:eq:right-unit-coherence}, stosujemy właśnie dowiedzioną tożsamość do kategorii z odwróconą kolejnością iloczynu: $A\odot B=B\otimes A$, $\alpha^\odot_{A,B,C}=\alpha^{-1}_{C,B,A}$, $\lambda^\odot_A=\rho_A$, $\rho^\odot_A=\lambda_A$. Pięciokąt dla $\odot$ jest odwróconym pięciokątem dla $\otimes$ z kolejnością czterech czynników odwróconą; trójkąt to równanie $(\rho_B\otimes\id_A)\alpha_{B,I,A}^{-1}=\id_B\otimes\lambda_A$, równoważne pierwotnemu. Pierwsza tożsamość w nowej strukturze brzmi $\rho_{B\otimes A}\alpha_{B,A,I}^{-1}=\id_B\otimes\rho_A$. Złożenie z $\alpha_{B,A,I}$ i zamiana nazw daje żądaną równość.

Naturalność $\lambda$ względem $\lambda_I:I\otimes I\to I$ daje $\lambda_I(\id_I\otimes\lambda_I)=\lambda_I\lambda_{I\otimes I}$. Skracamy izomorfizm $\lambda_I$, więc $\lambda_{I\otimes I}=\id_I\otimes\lambda_I$. Wstawiając do \eqref{cat:eq:left-unit-coherence} dla $A=B=I$ i stosując trójkąt, dostajemy $\lambda_I\otimes\id_I=\rho_I\otimes\id_I$. Wierność $-\otimes I$ daje \eqref{cat:eq:unitors-equal}.
\end{proof}

\subsection{Twierdzenie o koherencji}
Wyrażenia formalne budujemy z liter oznaczających obiekty, specjalnego symbolu jednostki $I$ oraz binarnego symbolu $\otimes$. Litera jest odróżniana od formalnego symbolu jednostki nawet wtedy, gdy przypisano jej obiekt jednostkowy. Słowo wyrażenia to uporządkowana lista liter po pominięciu formalnych jednostek.

\begin{twierdzenie}[Mac Lane'a o koherencji monoidalnej]\label{cat:thm:mac-lane-coherence}
Dla dwóch formalnych wyrażeń o tym samym słowie wszystkie morfizmy otrzymane z $\alpha,\lambda,\rho$, ich odwrotności, identyczności przez składanie i iloczyn monoidalny są równe. Istnieje co najmniej jeden taki izomorfizm. Porównujemy morfizmy o tych samych formalnych końcach, a nie dowolne morfizmy pomiędzy obiektami kategorii.
\end{twierdzenie}
\begin{proof}
Ustalamy orientację trzech operacji, które wolno wykonywać w dowolnym podwyrażeniu:
\[
((U\otimes V)\otimes W)\to U\otimes(V\otimes W),\qquad
I\otimes U\to U,\qquad U\otimes I\to U.
\]
Odpowiadają im odpowiednio $\alpha,\lambda,\rho$, z identycznościami w pozostałych czynnikach. Każdy krok jest izomorfizmem.

\textbf{Zakończenie redukcji.}
Niech $n(T)$ będzie liczbą wszystkich liści drzewa wyrażenia $T$ (łącznie z jednostkami), a $w(T)$ sumą, po wszystkich węzłach binarnych, liczby liści lewego poddrzewa. Porządkujemy pary $(n,w)$ leksykograficznie w $\mathbb N^2$. Usunięcie jednostki zmniejsza $n$. Zmiana nawiasowania nie zmienia $n$, ale zmniejsza $w$ o $n(U)$: wkład dwóch zmienionych węzłów wynosi przed krokiem $2n(U)+n(V)$, a po nim $n(U)+n(V)$. Wkłady węzłów poza podwyrażeniem pozostają takie same. Nie istnieje zatem nieskończony ciąg redukcji.

\textbf{Postać normalna.}
Każde wyrażenie nieredukowalne jest pojedynczą literą, samą jednostką albo wyrażeniem prawostronnie nawiasowanym złożonym z liter. Istotnie, gdy korzeń jest binarny, lewy czynnik nie może być binarny ani jednostką, prawy zaś nie może być jednostką. Stosujemy ten argument rekurencyjnie w prawym poddrzewie. Postać normalna $N(T)$ jest jednoznacznie wyznaczona przez słowo: dla pustego słowa jest to $I$. Każdy krok zachowuje słowo.

\textbf{Lokalne rozgałęzienia.}
Dwa różne kroki z tego samego wyrażenia można uzupełnić do diagramu przemiennego złożonego z redukcji. Jeżeli dotyczą rozłącznych poddrzew, wykonujemy oba w dowolnej kolejności; bifunktorialność iloczynu gwarantuje równość. Jeżeli jeden krok występuje wewnątrz zmiennej $U,V,W$ drugiej reguły, naturalność odpowiednio $\alpha,\lambda,\rho$ daje kwadrat przemienny. Pozostają nakładania w miejscach niebędących zmiennymi reguł. Są dokładnie następujące:
\[
\begin{array}{c|c}
\text{wyrażenie}&\text{równość uzupełniająca}\\ \hline
((U\otimes V)\otimes W)\otimes Z&\text{pięciokąt}\\
(I\otimes U)\otimes V&\eqref{cat:eq:left-unit-coherence}\\
(U\otimes I)\otimes V&\text{trójkąt}\\
(U\otimes V)\otimes I&\eqref{cat:eq:right-unit-coherence}\\
I\otimes I&\eqref{cat:eq:unitors-equal}
\end{array}
\]
Aby sprawdzić kompletność, zauważmy, że reguła $\alpha$ ma tylko jedno właściwe podwyrażenie binarne, które nie jest zmienną, mianowicie $U\otimes V$. Może ono rozpoczynać każdą z trzech reguł, co daje pierwsze trzy przypadki. Na korzeniu reguła $\alpha$ może zbiegać się z $\rho$, jeśli $W=I$, co daje czwarty przypadek. Reguły $\lambda,\rho$ mogą zbiegać się tylko w $I\otimes I$. Pozostałe właściwe podwyrażenia reguł jednostki są zmiennymi lub liśćmi. Jeżeli nakładanie występuje w większym kontekście, iloczyn z identycznościami zachowuje równość. Wszystkie wymienione diagramy są przemienne na mocy aksjomatów CAT-19 lub lematu~\ref{cat:lem:unit-coherence}.

\textbf{Globalna jednoznaczność.}
Indukujemy względem dobrze ufundowanej miary $(n,w)$. Dla wyrażenia normalnego jedyną drogą redukcji do postaci normalnej jest identyczność. Dla pozostałych wyrażeń rozważamy dwie drogi do postaci normalnej. Jeżeli zaczynają się tym samym krokiem, stosujemy założenie indukcyjne do jego celu. W przeciwnym razie oznaczamy pierwsze kroki przez $f:T\to U$, $g:T\to V$. Z lokalnej kontroli mamy drogi redukcji $p:U\to W$, $q:V\to W$ z $pf=qg$. Wybieramy drogę $r:W\to N(T)$, która istnieje dzięki zakończeniu redukcji i zachowaniu słowa. Miary $U,V$ są ściśle mniejsze od miary $T$ niezależnie od tego, czy $U=V$ jako wyrażenia formalne; także dwa różne kroki o wspólnym celu podlegają argumentowi lokalnemu (w szczególności $\lambda_I,\rho_I$). Wówczas $rp$ i $rq$ są drogami normalizującymi. Z założenia indukcyjnego są równe odpowiednim ogonom pierwotnych dróg. Równość $rpf=rqg$ dowodzi tezy.

Oznaczmy jedyny morfizm normalizujący przez $\nu_T:T\to N(T)$. Dla każdego elementarnego kroku $f:T\to U$ mamy $\nu_U f=\nu_T$, zatem $f=\nu_U^{-1}\nu_T$. Dla jego odwrotności zachodzi analogiczny wzór. Dowolny iloczyn dwóch morfizmów można rozpisać na kroki w podwyrażeniach, ponieważ $f\otimes g=(\id\otimes g)(f\otimes\id)$ z właściwymi identycznościami. W dowolnej drodze złożenia $\nu_U\nu_U^{-1}$ skracają się, pozostawiając $\nu_S^{-1}\nu_T$ dla końców $T,S$. Jest to również istniejący morfizm między tymi końcami. Dowód nie używa koherencji w przesłankach, tylko aksjomatów i lematu jednostki.
\end{proof}

\begin{uwaga}
Twierdzenie pozwala pomijać nawiasy w rachunkach strukturalnych, lecz nie utożsamia dowolnych morfizmów. Nie obejmuje samo z siebie permutacji czynników przez plecenie; dla kategorii symetrycznych udowodnimy odrębne twierdzenie o koherencji symetrii w ostatniej części rozdziału. Historyczny punkt odniesienia: S.~Mac Lane, \emph{Categories for the Working Mathematician}, rozdział VII; redukcja aksjomatów jednostki: G.~M.~Kelly, \emph{On MacLane's conditions for coherence of natural associativities, commutativities, etc.}, Journal of Algebra 1 (1964).
\end{uwaga}

\subsection{Funktory monoidalne}
\begin{definicja}[Funktor monoidalny typu lax]\label{cat:def:lax-monoidal}
Niech $\mathcal C,\mathcal D$ będą kategoriami monoidalnymi. Funktor monoidalny typu lax (o morfizmach strukturalnych skierowanych do obrazu iloczynu) składa się z funktora $F:\mathcal C\to\mathcal D$, naturalnej rodziny
\[
\phi_{A,B}:F(A)\otimes F(B)\to F(A\otimes B),\qquad
\phi_0:I_{\mathcal D}\to F(I_{\mathcal C})
\]
spełniających równości
\[
\begin{aligned}
&F(\alpha_{A,B,C})\phi_{A\otimes B,C}(\phi_{A,B}\otimes\id_{F(C)})\\
&=\phi_{A,B\otimes C}(\id_{F(A)}\otimes\phi_{B,C})\alpha_{F(A),F(B),F(C)},\\
&F(\lambda_A)\phi_{I_{\mathcal C},A}(\phi_0\otimes\id_{F(A)})=\lambda_{F(A)},\\
&F(\rho_A)\phi_{A,I_{\mathcal C}}(\id_{F(A)}\otimes\phi_0)=\rho_{F(A)}.
\end{aligned}
\]
W pierwszym równaniu dziedziną jest $(F(A)\otimes F(B))\otimes F(C)$, a przeciwdziedziną $F(A\otimes(B\otimes C))$; w dwóch ostatnich odpowiednio $I_{\mathcal D}\otimes F(A)$ oraz $F(A)\otimes I_{\mathcal D}$, z przeciwdziedziną $F(A)$. Funktor jest silny, jeśli wszystkie $\phi_{A,B}$ i $\phi_0$ są izomorfizmami; jest ścisły, jeśli są identycznościami (w szczególności wymagane są odpowiednie równości obiektów).
\end{definicja}

\begin{definicja}[Funktor komonoidalny]\label{cat:def:oplax-monoidal}
Funktor typu oplax (komonoidalny) ma naturalne morfizmy $\psi_{A,B}:F(A\otimes B)\to F(A)\otimes F(B)$ oraz $\psi_0:F(I_{\mathcal C})\to I_{\mathcal D}$, spełniające
\[
\begin{aligned}
&\alpha_{F(A),F(B),F(C)}(\psi_{A,B}\otimes\id_{F(C)})\psi_{A\otimes B,C}\\
&=(\id_{F(A)}\otimes\psi_{B,C})\psi_{A,B\otimes C}F(\alpha_{A,B,C}),\\
&\lambda_{F(A)}(\psi_0\otimes\id_{F(A)})\psi_{I_{\mathcal C},A}=F(\lambda_A),\\
&\rho_{F(A)}(\id_{F(A)}\otimes\psi_0)\psi_{A,I_{\mathcal C}}=F(\rho_A).
\end{aligned}
\]
Jeżeli struktura lax jest silna, jej odwrotne izomorfizmy definiują strukturę oplax: odwracamy obie drogi w każdym diagramie i w razie potrzeby składamy z izomorfizmem łączności lub jednostki, aby uzyskać powyższe równania.
\end{definicja}

\begin{definicja}[Monoidalna transformacja naturalna]\label{cat:def:monoidal-natural}
Dla funktorów lax $F,G$ o strukturach $\phi,\gamma$ transformacja naturalna $\eta:F\Rightarrow G$ jest monoidalna, jeśli
\[
\eta_{A\otimes B}\phi_{A,B}
=\gamma_{A,B}(\eta_A\otimes\eta_B),\qquad
\eta_{I_{\mathcal C}}\phi_0=\gamma_0.
\]
Dla funktorów silnych równoważne warunki otrzymujemy używając odwrotnych morfizmów strukturalnych.
\end{definicja}

\begin{stwierdzenie}[Składanie funktorów monoidalnych]\label{cat:prop:monoidal-composition}
Niech $F:\mathcal C\to\mathcal D$ oraz $G:\mathcal D\to\mathcal E$ będą funktorami typu lax o strukturach odpowiednio $\phi$ i $\gamma$. Wówczas $GF$ ma strukturę typu lax
\[
\theta_{A,B}=G(\phi_{A,B})\gamma_{F(A),F(B)},\qquad
\theta_0=G(\phi_0)\gamma_0.
\]
Jeśli oba funktory są silne (odpowiednio ścisłe), to ich złożenie również jest silne (odpowiednio ścisłe).
\end{stwierdzenie}
\begin{proof}
Typy morfizmów są następujące:
\[
GF(A)\otimes GF(B)\xrightarrow{\gamma_{F(A),F(B)}}
G(F(A)\otimes F(B))\xrightarrow{G(\phi_{A,B})}GF(A\otimes B).
\]
Morfizm jednostki prowadzi z $I_{\mathcal E}$ przez $G(I_{\mathcal D})$ do $GF(I_{\mathcal C})$. Naturalność $\theta$ wynika z naturalności $\phi,\gamma$: dla $a:A\to A'$, $b:B\to B'$ mamy
\[
\begin{aligned}
GF(a\otimes b)\theta_{A,B}
&=G(\phi_{A',B'})G(F(a)\otimes F(b))\gamma_{F(A),F(B)}\\
&=G(\phi_{A',B'})\gamma_{F(A'),F(B')}(GF(a)\otimes GF(b)).
\end{aligned}
\]
W rachunku łączności zapisujemy $X=F(A),Y=F(B),Z=F(C)$ oraz $P=F(A\otimes B)$, $Q=F(B\otimes C)$. Wstawiając definicję $\theta$ i stosując naturalność $\gamma$ względem $\phi_{A,B}:X\otimes Y\to P$, otrzymujemy
\[
\begin{aligned}
&GF(\alpha_{A,B,C})\theta_{A\otimes B,C}(\theta_{A,B}\otimes\id)\\
&=G(F(\alpha_{A,B,C})\phi_{A\otimes B,C}(\phi_{A,B}\otimes\id_Z))\\
&\qquad{}\cdot\gamma_{X\otimes Y,Z}(\gamma_{X,Y}\otimes\id_{G(Z)})\\
&=G(\phi_{A,B\otimes C}(\id_X\otimes\phi_{B,C}))\\
&\qquad{}\cdot G(\alpha_{X,Y,Z})\gamma_{X\otimes Y,Z}(\gamma_{X,Y}\otimes\id_{G(Z)})\\
&=G(\phi_{A,B\otimes C})G(\id_X\otimes\phi_{B,C})\\
&\qquad{}\cdot\gamma_{X,Y\otimes Z}(\id_{G(X)}\otimes\gamma_{Y,Z})
\alpha_{G(X),G(Y),G(Z)}\\
&=G(\phi_{A,B\otimes C})\gamma_{X,Q}
(\id_{G(X)}\otimes G(\phi_{B,C}))\\
&\qquad{}\cdot(\id_{G(X)}\otimes\gamma_{Y,Z})\alpha_{G(X),G(Y),G(Z)}\\
&=\theta_{A,B\otimes C}(\id_{GF(A)}\otimes\theta_{B,C})
\alpha_{GF(A),GF(B),GF(C)}.
\end{aligned}
\]
Kolejne równości wykorzystują naturalność $\gamma$, aksjomat łączności $F$, aksjomat łączności $G$, naturalność $\gamma$ i bifunktorialność. Wszystkie złożenia mają wspólne końce $(GF(A)\otimes GF(B))\otimes GF(C)$ i $GF(A\otimes(B\otimes C))$.

Dla lewej jednostki naturalność $\gamma$ względem $\phi_0$ daje
\[
\begin{aligned}
GF(\lambda_A)\theta_{I_{\mathcal C},A}(\theta_0\otimes\id)
&=G(F(\lambda_A)\phi_{I_{\mathcal C},A}(\phi_0\otimes\id_X))\\
&\quad{}\cdot\gamma_{I_{\mathcal D},X}(\gamma_0\otimes\id_{G(X)})\\
&=G(\lambda_X)\gamma_{I_{\mathcal D},X}(\gamma_0\otimes\id_{G(X)})\\
&=\lambda_{G(X)}.
\end{aligned}
\]
Analogiczny rachunek dla prawej jednostki, z podaniem odpowiednich czynników, to
\[
\begin{aligned}
GF(\rho_A)\theta_{A,I_{\mathcal C}}(\id\otimes\theta_0)
&=G(F(\rho_A)\phi_{A,I_{\mathcal C}}(\id_X\otimes\phi_0))\\
&\quad{}\cdot\gamma_{X,I_{\mathcal D}}(\id_{G(X)}\otimes\gamma_0)\\
&=G(\rho_X)\gamma_{X,I_{\mathcal D}}(\id_{G(X)}\otimes\gamma_0)\\
&=\rho_{G(X)}.
\end{aligned}
\]
Funktory zachowują izomorfizmy i identyczności, więc złożenia struktur silnych lub ścisłych mają tę samą własność.
\end{proof}

\begin{stwierdzenie}[Składanie monoidalnych transformacji]\label{cat:prop:monoidal-transform-composition}
Identyczności oraz złożenia pionowe i poziome monoidalnych transformacji naturalnych między funktorami typu lax są monoidalne, przy strukturach funktorów złożonych z poprzedniego stwierdzenia.
\end{stwierdzenie}
\begin{proof}
Dla identyczności warunki są bezpośrednimi równościami struktur. Niech $\eta:F\Rightarrow G$, $\zeta:G\Rightarrow H$ mają struktury odpowiednio $\phi,\gamma,\chi$. Wtedy
\[
(\zeta\eta)_{A\otimes B}\phi_{A,B}
=\zeta_{A\otimes B}\gamma_{A,B}(\eta_A\otimes\eta_B)
=\chi_{A,B}((\zeta_A\eta_A)\otimes(\zeta_B\eta_B)),
\]
a $(\zeta\eta)_I\phi_0=\zeta_I\gamma_0=\chi_0$.
Dla składania poziomego wystarczy sprawdzić oba rodzaje składania z identycznością funktora, ponieważ dowolne złożenie poziome jest ich złożeniem pionowym.
Jeżeli $K:\mathcal D\to\mathcal E$ ma strukturę $\kappa$, to
\[
\begin{aligned}
K(\eta_{A\otimes B})K(\phi_{A,B})\kappa_{F(A),F(B)}
&=K(\gamma_{A,B})K(\eta_A\otimes\eta_B)\kappa_{F(A),F(B)}\\
&=K(\gamma_{A,B})\kappa_{G(A),G(B)}(K(\eta_A)\otimes K(\eta_B)).
\end{aligned}
\]
Jednostka spełnia $K(\eta_I)K(\phi_0)\kappa_0=K(\gamma_0)\kappa_0$.
Jeżeli $L:\mathcal B\to\mathcal C$ ma strukturę $\ell$, to naturalność $\eta$ względem $\ell_{A,B}:L(A)\otimes L(B)\to L(A\otimes B)$ daje
\[
\begin{aligned}
\eta_{L(A\otimes B)}F(\ell_{A,B})\phi_{L(A),L(B)}
&=G(\ell_{A,B})\eta_{L(A)\otimes L(B)}\phi_{L(A),L(B)}\\
&=G(\ell_{A,B})\gamma_{L(A),L(B)}(\eta_{L(A)}\otimes\eta_{L(B)}).
\end{aligned}
\]
Podobnie $\eta_{L(I)}F(\ell_0)\phi_0=G(\ell_0)\eta_{I_{\mathcal C}}\phi_0=G(\ell_0)\gamma_0$.
Otrzymaliśmy wszystkie wymagane warunki.
\end{proof}

\begin{definicja}[Równoważność monoidalna]\label{cat:def:monoidal-equivalence}
Równoważnością monoidalną nazywamy parę silnych funktorów monoidalnych w przeciwnych kierunkach wraz z monoidalnymi izomorfizmami naturalnymi ich złożeń z odpowiednimi funktorami identycznościowymi. Funktory identycznościowe mają ścisłą strukturę monoidalną.
\end{definicja}

\subsection{Uściślanie kategorii monoidalnej}
\begin{twierdzenie}[Uściślanie Mac Lane'a]\label{cat:thm:strictification}
Każda lokalnie mała kategoria monoidalna jest silnie monoidalnie równoważna pewnej ścisłej kategorii monoidalnej, przy ustalonej konwencji wielkości.
\end{twierdzenie}
\begin{proof}
Niech obiektami $\mathcal S$ będą skończone listy obiektów $\mathcal C$, także lista pusta. Dla listy $u$ definiujemy $E(u)$ jako prawostronnie nawiasowany iloczyn jej elementów, przyjmując $E(())=I$ oraz $E((A))=A$. Definiujemy $\Hom_{\mathcal S}(u,v)=\Hom_{\mathcal C}(E(u),E(v))$ (morfizmy mają oznaczone końce $u,v$), składanie i identyczności odziedziczone z $\mathcal C$. Aksjomaty kategorii wynikają z aksjomatów $\mathcal C$.

Iloczyn obiektów to konkatenacja $uv$, jednostka to lista pusta. Z twierdzenia~\ref{cat:thm:mac-lane-coherence} otrzymujemy kanoniczne izomorfizmy $m_{u,v}:E(u)\otimes E(v)\to E(uv)$, stosując formalne litery dla pozycji listy. Dla morfizmów $f:u\to u'$, $g:v\to v'$ definiujemy
\[
f\otimes_{\mathcal S}g=m_{u',v'}(f\otimes g)m_{u,v}^{-1}.
\]
Jest to morfizm $uv\to u'v'$. Zachowanie identyczności wynika z $m_{u,v}m_{u,v}^{-1}=\id$. Dla dalszych morfizmów $f',g'$ środkowe izomorfizmy $m_{u',v'}^{-1}m_{u',v'}$ skracają się; bifunktorialność w $\mathcal C$ daje $(f'\otimes_{\mathcal S}g')(f\otimes_{\mathcal S}g)=(f'f)\otimes_{\mathcal S}(g'g)$.

Oznaczmy
\[
L_{u,v,w}=m_{uv,w}(m_{u,v}\otimes\id),\qquad
R_{u,v,w}=m_{u,vw}(\id\otimes m_{v,w}).
\]
Koherencja daje $L_{u,v,w}=R_{u,v,w}\alpha_{E(u),E(v),E(w)}$. Dla $h:w\to w'$ rozwinięcie definicji daje
\[
\begin{aligned}
(f\otimes_{\mathcal S}g)\otimes_{\mathcal S}h
&=L_{u',v',w'}((f\otimes g)\otimes h)L_{u,v,w}^{-1},\\
f\otimes_{\mathcal S}(g\otimes_{\mathcal S}h)
&=R_{u',v',w'}(f\otimes(g\otimes h))R_{u,v,w}^{-1}.
\end{aligned}
\]
Po podstawieniu $L=R\alpha$ pierwsze wyrażenie redukuje się do drugiego przez naturalność $\alpha$. Ponadto $m_{(),u}=\lambda_{E(u)}$, $m_{u,()}=\rho_{E(u)}$ z koherencji; naturalność jednostek daje $\id_{()}\otimes_{\mathcal S}f=f=f\otimes_{\mathcal S}\id_{()}$. Wobec tego $\mathcal S$ jest ściśle monoidalna.

Funktor $E:\mathcal S\to\mathcal C$ jest identycznością na odpowiednich zbiorach morfizmów, zatem jest w pełni wierny. Każdy $A$ jest obrazem listy $(A)$. Struktura silna to $m_{u,v}$ i $\id_I$. Naturalność $m$ wynika z definicji iloczynu morfizmów, aksjomat łączności to $L=R\alpha$, a aksjomaty jednostek wynikają z podanych wzorów dla pustej listy.

Jawny quasi-odwrotny funktor $J:\mathcal C\to\mathcal S$ wysyła $A$ na $(A)$ i $f$ na ten sam morfizm w odpowiednim zbiorze. Jego struktura $J(A)J(B)\to J(A\otimes B)$ oraz $()\to J(I)$ jest reprezentowana w $\mathcal C$ przez odpowiednie identyczności. Sprawdźmy jawnie aksjomaty $J$. Oznaczmy jego porównania przez $j_{A,B}$ i $j_0$. Z definicji $m_{(A),(B)}=\id_{A\otimes B}$. Naturalność $j$ po zastosowaniu wiernego $E$ sprowadza się więc do równości $a\otimes b=a\otimes b$ dla morfizmów $a,b$ w $\mathcal C$.

Dla list $u=(A),v=(B),w=(C)$ mamy $m_{uv,w}=\alpha_{A,B,C}$ i $m_{u,vw}=\id_{A\otimes(B\otimes C)}$; pozostałe porównania pojedynczych list są identycznościami. Dlatego
\[
\begin{aligned}
E(j_{A,B}\otimes_{\mathcal S}\id_{J(C)})&=\alpha_{A,B,C}^{-1},\\
E(\id_{J(A)}\otimes_{\mathcal S}j_{B,C})&=\id_{A\otimes(B\otimes C)}.
\end{aligned}
\]
Po zastosowaniu $E$ lewa droga aksjomatu łączności $J$ daje $\alpha_{A,B,C}\alpha_{A,B,C}^{-1}$, prawa identyczność (łączność $\mathcal S$ jest ścisła). Dla jednostki $E(j_0\otimes_{\mathcal S}\id)=\lambda_A^{-1}$, gdyż $m_{(),(A)}=\lambda_A$ i $m_{(I),(A)}=\id$. Podobnie $E(\id\otimes_{\mathcal S}j_0)=\rho_A^{-1}$. Złożenia z $E(J(\lambda_A))$ i $E(J(\rho_A))$ dają identyczności. Wierność $E$ dowodzi wszystkich aksjomatów; struktura $J$ jest silna, ponieważ jej składowe są reprezentowane przez identyczności. Mamy $EJ=\id_{\mathcal C}$ jako funktory silne: według stwierdzenia~\ref{cat:prop:monoidal-composition} struktura złożenia to $E(j_{A,B})m_{(A),(B)}=\id$ oraz $E(j_0)\id_I=\id_I$. Izomorfizm $JE\Rightarrow\id_{\mathcal S}$ w $u$ jest reprezentowany przez $\id_{E(u)}$, jest naturalny z definicji składania. Warunek monoidalności dla pary list po zastosowaniu $E$ ma po obu stronach morfizm $m_{u,v}:E(u)\otimes E(v)\to E(uv)$; warunek jednostki ma po obu stronach $\id_I$. Wierność $E$ daje wymagane równości. Otrzymaliśmy monoidalną równoważność.
\end{proof}

\begin{stwierdzenie}[Porównania kartezjańskie]\label{cat:prop:cartesian-oplax}
Każdy funktor między kategoriami z wybranymi skończonymi iloczynami, wyposażonymi w struktury kartezjańskie, ma kanoniczną strukturę oplax. Jest ona odwracalna dokładnie wtedy, gdy funktor zachowuje skończone iloczyny.
\end{stwierdzenie}
\begin{proof}
Definiujemy $\psi_{A,B}=\langle F(p_A),F(p_B)\rangle$ oraz $\psi_0$ jako jedyny morfizm do obiektu końcowego. Naturalność $\psi$ wynika z porównania obu rzutów i funktorialności $F$. Obie drogi w aksjomacie łączności mają składowe $F(p_A),F(p_B),F(p_C)$ dla złożonych rzutów z $(A\times B)\times C$, ponieważ kanoniczny izomorfizm łączności zachowuje te rzuty. Jednoznaczność morfizmu do iterowanego iloczynu daje równość. W aksjomatach jednostki po złożeniu z odpowiednim rzutem otrzymujemy $F(\lambda_A)$ lub $F(\rho_A)$, co daje wymaganą równość. Porównanie $\psi_{A,B}$ jest izomorfizmem wtedy i tylko wtedy, gdy $F(A\times B)$ z rzutami $F(p_A),F(p_B)$ jest iloczynem. Analogicznie $\psi_0$ jest izomorfizmem wtedy i tylko wtedy, gdy $F(1)$ jest końcowy. Skończone iloczyny otrzymujemy przez iterację binarnych oraz pusty iloczyn. To kończy dowód.
\end{proof}

\begin{definicja}[Zachowywanie plecenia]\label{cat:def:braided-functor}
Silny funktor monoidalny między kategoriami z pleceniem zachowuje plecenie, jeśli
\[
F(c^{\mathcal C}_{A,B})\phi_{A,B}
=\phi_{B,A}c^{\mathcal D}_{F(A),F(B)}.
\]
Obie strony są morfizmami $F(A)\otimes F(B)\to F(B\otimes A)$. Gdy obie kategorie są symetryczne, taki funktor nazywamy symetrycznym monoidalnym.
\end{definicja}

\begin{przyklad}
Funktor zachowujący skończone iloczyny z poprzedniego stwierdzenia, wyposażony w $\phi=\psi^{-1}$, jest symetryczny monoidalny: równoważna równość $\psi_{B,A}F(c_{A,B})=c_{F(A),F(B)}\psi_{A,B}$ zachodzi, ponieważ obie strony mają składowe $F(p_B),F(p_A)$.
\end{przyklad}

\subsection{Koherencja symetryczna}
\begin{lemat}[Przeniesienie symetrii do kategorii ścisłej]\label{cat:lem:strict-symmetry}
Jeżeli $\mathcal C$ jest symetryczna monoidalna, to jej uściślenie $\mathcal S$ z twierdzenia~\ref{cat:thm:strictification} jest ścisłe symetryczne monoidalne dla
\[
E(c^{\mathcal S}_{u,v})=m_{v,u}c_{E(u),E(v)}m_{u,v}^{-1}.
\]
Funktor $E$ zachowuje symetrię.
\end{lemat}
\begin{proof}
Prawa strona ma typ $E(uv)\to E(vu)$, więc określa dokładnie jeden morfizm w $\mathcal S$. Naturalność po podstawieniu wzoru na iloczyn morfizmów redukuje się do naturalności $c$ w $\mathcal C$: dla $f:u\to u'$, $g:v\to v'$ obie strony, po zastosowaniu $E$, są równe
\[
m_{v',u'}(g\otimes f)c_{E(u),E(v)}m_{u,v}^{-1}.
\]
Podwójna zamiana daje $m_{u,v}c_{E(v),E(u)}c_{E(u),E(v)}m_{u,v}^{-1}=\id$. Warunek zachowania symetrii przez $E$ jest dokładnie definicją $c^{\mathcal S}$.

Sprawdzimy sześciokąty. Zapiszmy $X=E(u)$, $Y=E(v)$, $Z=E(w)$ i użyjmy porównań $L_{u,v,w}$ z dowodu uściślania, identyfikujących lewostronnie nawiasowane trójki z wartościami konkatenacji. Po zastosowaniu $E$ i sprzężeniu morfizmów tymi porównaniami, tj. przez złożenie po prawej z $L$ dla źródła i po lewej z $L^{-1}$ dla celu, dostajemy dla pierwszego ścisłego sześciokąta:
\[
\begin{aligned}
c^{\mathcal S}_{u,vw}&\ \longmapsto\
 c_{X,Y\otimes Z}\alpha_{X,Y,Z},\\
(\id_v\otimes_{\mathcal S} c^{\mathcal S}_{u,w})
(c^{\mathcal S}_{u,v}\otimes_{\mathcal S}\id_w)
&\ \longmapsto\
 \alpha_{Y,Z,X}^{-1}(\id_Y\otimes c_{X,Z})\\
&\qquad{}\cdot\alpha_{Y,X,Z}(c_{X,Y}\otimes\id_Z).
\end{aligned}
\]
Istotnie, dla pierwszego morfizmu $L_{v,w,u}^{-1}m_{vw,u}=m_{v,w}^{-1}\otimes\id_X$ i $m_{u,vw}^{-1}L_{u,v,w}=(\id_X\otimes m_{v,w})\alpha_{X,Y,Z}$. Naturalność $c$ skraca $m_{v,w}$ z jego odwrotnością. W drugiej drodze pierwszy krok sprzęgamy przez $L_{v,u,w}^{-1}$ i $L_{u,v,w}$, co daje $c_{X,Y}\otimes\id_Z$. Dla drugiego kroku przechodzimy do prawostronnego nawiasowania przez $\alpha_{Y,X,Z}$, stosujemy $\id_Y\otimes c_{X,Z}$ i wracamy przez $\alpha_{Y,Z,X}^{-1}$; wynika to ze wzoru $L=R\alpha$. Równość otrzymanych wyrażeń jest pierwszym sześciokątem z CAT-19 po złożeniu z $\alpha_{Y,Z,X}^{-1}$.

Dla drugiego ścisłego sześciokąta ten sam wzór $L=R\alpha$ daje
\[
\begin{aligned}
c^{\mathcal S}_{uv,w}&\ \longmapsto\
 \alpha_{Z,X,Y}^{-1}c_{X\otimes Y,Z},\\
(c^{\mathcal S}_{u,w}\otimes_{\mathcal S}\id_v)
(\id_u\otimes_{\mathcal S} c^{\mathcal S}_{v,w})
&\ \longmapsto\
(c_{X,Z}\otimes\id_Y)\alpha_{X,Z,Y}^{-1}\\
&\qquad{}\cdot(\id_X\otimes c_{Y,Z})\alpha_{X,Y,Z}.
\end{aligned}
\]
W pierwszej drodze $L_{w,u,v}^{-1}m_{w,uv}=\alpha_{Z,X,Y}^{-1}(\id_Z\otimes m_{u,v}^{-1})$ i $m_{uv,w}^{-1}L_{u,v,w}=m_{u,v}\otimes\id_Z$, więc naturalność skraca porównania. W drugiej drodze pierwszy krok wymaga zmiany nawiasowania, a drugi działa na dwóch pierwszych czynnikach. Drugi sześciokąt z CAT-19, złożony po prawej z $\alpha_{X,Y,Z}$, daje równość. W obu przypadkach $E$ jest wierny, a porównania odwracalne, zatem ścisłe sześciokąty zachodzą również w $\mathcal S$. Rachunki obejmują puste listy, ponieważ wzór $L=R\alpha$ został dowiedziony także dla nich.
\end{proof}

\begin{lemat}[Relacje sąsiednich zamian]\label{cat:lem:coxeter-swaps}
W ścisłej symetrycznej kategorii monoidalnej zamiana czynników na pozycjach $i,i+1$, dana przez $c$ i identyczności na pozostałych pozycjach, spełnia relacje
\[
s_i^2=1,\qquad s_i s_j=s_j s_i\ (|i-j|>1),\qquad
s_i s_{i+1}s_i=s_{i+1}s_i s_{i+1}.
\]
W każdym kroku $s_i$ oznacza zamianę aktualnych czynników, a nie ustalony endomorfizm niezależny od słowa.
\end{lemat}
\begin{proof}
Pierwsza relacja jest symetrią, druga wynika z bifunktorialności dla rozłącznych par pozycji. Trzecia sprowadza się do trzech kolejnych czynników $A,B,C$. Naturalność $c$ względem morfizmu $c_{A,B}:A\otimes B\to B\otimes A$ i $\id_C$ daje
\[
c_{B\otimes A,C}(c_{A,B}\otimes\id_C)
=(\id_C\otimes c_{A,B})c_{A\otimes B,C}.
\]
Rozwinięcie obu zamian blokowych za pomocą drugiego ścisłego sześciokąta daje
\[
\begin{aligned}
&(c_{B,C}\otimes\id_A)(\id_B\otimes c_{A,C})(c_{A,B}\otimes\id_C)\\
&=(\id_C\otimes c_{A,B})(c_{A,C}\otimes\id_B)(\id_A\otimes c_{B,C}).
\end{aligned}
\]
To dokładnie trzecia relacja; obie drogi prowadzą z $A\otimes B\otimes C$ do $C\otimes B\otimes A$. Iloczyn z identycznościami przenosi tę równość na dowolne trzy kolejne pozycje.
\end{proof}

\begin{lemat}[Prezentacja grupy permutacji]\label{cat:lem:symmetric-presentation}
Dwa słowa w sąsiednich transpozycjach $s_1,\ldots,s_{n-1}$ wyznaczają tę samą permutację $n$ pozycji wtedy i tylko wtedy, gdy można przejść od jednego do drugiego przez relacje z lematu~\ref{cat:lem:coxeter-swaps}.
\end{lemat}
\begin{proof}
Dla $n=0,1$ twierdzenie jest bezpośrednie. Rozważmy grupę $G_n$ określoną przez podane generatory i relacje. Ponieważ $s_i^2=1$, odwrotności generatorów można zastąpić generatorami. Istnieje surjektywny homomorfizm $G_n\to S_n$, gdyż sąsiednie transpozycje spełniają relacje i generują $S_n$ (dowolną literę można przesunąć na żądane miejsce przez kolejne zamiany sąsiadów).

Niech $H$ będzie podgrupą generowaną przez $s_1,\ldots,s_{n-2}$. Jest obrazem $G_{n-1}$, zatem z indukcji ma co najwyżej $(n-1)!$ elementów. Definiujemy $r_j=s_{n-1}s_{n-2}\cdots s_j$ dla $1\le j<n$ oraz $r_n=1$. Pokażemy, że każdy element ma postać $hr_j$, $h\in H$. Zbiór tych postaci zawiera identyczność i jest zamknięty na mnożenie z prawej przez dowolny $s_i$. Dla $j<n$ mamy
\[
r_js_i=
\begin{cases}
s_i r_j,&i<j-1,\\
r_{j-1},&i=j-1,\\
r_{j+1},&i=j,\\
s_{i-1}r_j,&i>j.
\end{cases}
\]
Pierwszy przypadek wynika z przemienności rozłącznych zamian, drugi z definicji, trzeci z $s_j^2=1$. W ostatnim przypadku przenosimy prawy $s_i$ przez czynniki $s_j,\ldots,s_{i-2}$, otrzymując fragment $s_i s_{i-1}s_i$, zamieniamy go na $s_{i-1}s_i s_{i-1}$ i przenosimy pierwszy $s_{i-1}$ w lewo przez $s_{i+1},\ldots,s_{n-1}$. Dla $j=n$ mamy $r_ns_i=s_i r_n$ przy $i<n-1$ oraz $r_ns_{n-1}=r_{n-1}$. Wszystkie dodatkowe generatory po lewej należą do $H$. Indukcja po długości słowa dowodzi żądanej postaci.
Stąd $|G_n|\le n|H|\le n!$. Surjekcja na $S_n$ wymusza równość i jest izomorfizmem. Równość w grupie z prezentacji oznacza dokładnie możliwość stosowania relacji (w tym w dowolnym kontekście słowa), co daje tezę.
\end{proof}

\begin{twierdzenie}[Koherencja symetryczna]\label{cat:thm:symmetric-coherence}
W lokalnie małej symetrycznej kategorii monoidalnej dwa morfizmy strukturalne o tych samych formalnych końcach, zbudowane z $\alpha,\lambda,\rho,c$, ich odwrotności, identyczności, składania i iloczynu monoidalnego, są równe, jeśli wyznaczają tę samą permutację pozycji liter po pominięciu formalnych jednostek. Każda permutacja zgodna z literami źródła i celu jest realizowalna przez taki morfizm.
\end{twierdzenie}
\begin{proof}
Najpierw pracujemy w kategorii ścisłej. Ścisły pierwszy sześciokąt z dwoma czynnikami jednostkowymi daje $c_{A,I}=c_{A,I}c_{A,I}$, więc z odwracalności $c_{A,I}=\id_A$. Drugi sześciokąt daje analogicznie $c_{I,A}=\id_A$. Zamiana pustego bloku jest zatem identycznością.
Każdą zamianę dwóch niepustych bloków można rozłożyć na zamiany sąsiednich pojedynczych czynników: pierwszy sześciokąt rozkłada zamianę z iloczynem w drugim argumencie, a drugi rozkłada zamianę z iloczynem w pierwszym. Indukcja względem łącznej liczby czynników obu bloków (z przypadkiem pustym i parą jednoelementową jako bazami) daje rozkład. Odwrotność zamiany jest zamianą w przeciwnym kierunku przez symetrię. Iloczyn morfizmów rozkładamy na działanie kolejno w dwóch blokach przez bifunktorialność. W ścisłej kategorii pozostałe morfizmy strukturalne są identycznościami. Otrzymujemy słowo zamian sąsiadów, realizujące dokładnie permutację pierwotnego morfizmu. Dwa słowa o tej samej permutacji są połączone relacjami z lematu~\ref{cat:lem:symmetric-presentation}. Każda relacja daje równość morfizmów dla aktualnej kolejności czynników na mocy lematu~\ref{cat:lem:coxeter-swaps}. Złożenie z wcześniejszymi i późniejszymi zamianami zachowuje równość. Każdą permutację można uzyskać przez zamiany sąsiadów.

Dla ogólnej $\mathcal C$ stosujemy uściślenie $\mathcal S$ i lemat~\ref{cat:lem:strict-symmetry}. Dla formalnego wyrażenia $T$ oznaczmy przez $u_T$ listę obiektów odpowiadających literom i przez $\nu_T:T\to E(u_T)$ kanoniczny izomorfizm z twierdzenia~\ref{cat:thm:mac-lane-coherence}. Dla morfizmu strukturalnego $f:T\to T'$ przenosimy go do $\mathcal S$ jako morfizm reprezentowany przez $\nu_{T'}f\nu_T^{-1}$. Składanie przenosi się bezpośrednio przez skrócenie $\nu$ z odwrotnością. Z koherencji monoidalnej
\[
\nu_{T\otimes U}=m_{u_T,u_U}(\nu_T\otimes\nu_U).
\]
Dlatego przeniesienie iloczynu dwóch morfizmów jest iloczynem ich przeniesień w $\mathcal S$. Przeniesienia $\alpha,\lambda,\rho$ oraz ich odwrotności są identycznościami, gdyż normalizacja nie zależy od drogi. Dla zamiany $c_{T,U}$ naturalność $c$ i powyższa równość dają
\[
\nu_{U\otimes T}c_{T,U}\nu_{T\otimes U}^{-1}
=m_{u_U,u_T}c_{E(u_T),E(u_U)}m_{u_T,u_U}^{-1},
\]
czyli dokładnie reprezentanta $c^{\mathcal S}_{u_T,u_U}$. Indukcja po budowie morfizmu pokazuje, że jego przeniesienie jest morfizmem strukturalnym w ścisłej kategorii z tą samą permutacją. Tam wykazaliśmy jednoznaczność. Ponieważ $\nu_T,\nu_{T'}$ są izomorfizmami, równość przeniesień daje równość pierwotnych morfizmów. Istnienie otrzymujemy przez normalizację, realizację sąsiednich zamian z kanonicznymi zmianami nawiasowania i odwrotną normalizację. Gdy litery oznaczają ten sam obiekt, nadal rozróżniamy ich pozycje; sama równość słów obiektów nie wystarcza do stwierdzenia równości morfizmów.
\end{proof}
